% 【本文件写什么】impl 实现层：virtio-pci
% - 定位：VirtIO PCI 传输，LoongArch 主线块设备。
% - crate 路径：os/components/wateros-driver/driver-block/block-impl/impl-virtio-pci/src/
% - 如何实现对应 api-v0 trait：关键类型、状态机、数据结构
% - 算法或硬件交互流程，可用伪代码/流程图
% - 本 impl 启用的 Cargo [features] 与 arch feature，impl-riscv64 等
% - 与其它 impl 变体的差异、切换 feature 时的影响
% - 性能/内存假设与已知限制
% 【事实来源】
% - 上述 crate 源码与 Cargo.toml
% - os/feature-tree.txt 中指向本 impl 的 feature 链

\subsubsection{impl-virtio-pci}

\paragraph{启用链。}由\code{impl-qemu-loongarch64-virt}打开。平台\code{pci::probe\_virtio\_blk\_pci}在PCIe ECAM bus0上枚举\code{virtio\_device\_type==Block}的首个功能，再经\code{VirtioPciBlkDevice::from\_pci\_root}完成BAR编程与传输握手。

\paragraph{BAR分配。}裸机bring-up无固件分配BAR时，\code{VirtioPciBarAllocator}在\code{[start,end)}区间内按\code{size}的2的幂对齐单调分配MMIO BAR；\code{Below1MiB}类型拒绝。\code{assign\_memory\_bars}遍历配置空间，对32/64位memory BAR调用\code{set\_bar\_32}/\code{set\_bar\_64}，并开启\code{MEMORY\_SPACE|BUS\_MASTER}。

\paragraph{探测入口。}\code{probe\_first\_from\_ecam}、\code{probe\_first\_from\_mmio\_cam}与\code{probe\_first\_from\_config}共享\code{probe\_first\_from\_config}实现，区别仅在\code{Cam::Ecam}与\code{Cam::MmioCam}。返回\code{(VirtioPciBlkDevice, VirtioPciProbeInfo)}供平台impl记录BDF与vendor/device id。

\begin{rustcode}
pub struct VirtioPciBlkDevice {
    inner: VirtIOBlk<VirtioPciHal, PciTransport>,
}

pub unsafe fn probe_first_from_ecam(
    ecam_base: usize,
    bar_allocator: &mut VirtioPciBarAllocator,
) -> DriverResult<Option<(Self, VirtioPciProbeInfo)>> {
    unsafe { Self::probe_first_from_config(ecam_base, Cam::Ecam, bar_allocator) }
}
\end{rustcode}

\paragraph{DMA与限制。}\code{VirtioPciHal}与MMIO路径相同：栈式帧分配器返回递减连续PPN，失败时释放已拿页并返回空指针。当前仅扫描bus0首个块设备；多盘场景需扩展枚举策略。块读写语义与MMIO实现一致，错误映射为\code{IoError}。
